\chapter{Notation and Symbol Table}
\label{app:notation}

\begin{longtable}{p{0.22\textwidth}p{0.68\textwidth}}
\toprule
Symbol & Meaning\\
\midrule
\endfirsthead
\toprule
Symbol & Meaning\\
\midrule
\endhead
$s=\sigma+it$ & Complex zeta variable, with real part $\sigma$ and imaginary part $t$.\\
$\Sstrip$ & Open critical strip $\{0<\Re s<1\}$.\\
$\Crit$ & Critical line $\{\Re s=1/2\}$.\\
$\zeta(s)$ & Riemann zeta function.\\
$\eta(s)$ & Dirichlet eta function $\sum(-1)^{n-1}n^{-s}$.\\
$\xi(s)$ & Completed zeta function $\frac12s(s-1)\pi^{-s/2}\Gamma(s/2)\zeta(s)$.\\
$\XiR(z)$ & Centred entire function $\xi(1/2+iz)$ in the early spectral convention.\\
$\J(s)$ & Conjugate-affine, anti-holomorphic involution $1-\bar s$.\\
$X$ & Channel swap operator on $\C^2$.\\
$\Hbin(\sigma)$ & Binary Shannon entropy in nats.\\
$\Acomp(\sigma,\phi)$ & Equal-gain complementary amplitude $\sqrt\sigma+e^{i\phi}\sqrt{1-\sigma}$.\\
$\ket{+},\ket{-}$ & Symmetric and antisymmetric detector rays.\\
$P_+$ & Projector $\ket{+}\bra{+}$.\\
$x$ & Horizontal displacement $\sigma-1/2$ in the early half-axis chapters.\\
$z$ & A centred complex coordinate; the local convention is stated in each chapter because both $z=s-1/2$ and $z=-i(s-1/2)$ occur in distinct reductions.\\
$u=\Omega(s)$ & Projective coordinate $s/(1-s)$ in the Li/inverse-boundary line. In G024, $u$ is redeclared locally as the non-negative centred correlation coordinate with $q=u^2$.\\
$q$ & Chapter-local compact or radial-square coordinate. In G024, $q=u^2\ge0$ and $H_y(q)=D_y(\sqrt q)$.\\
$w$ & Even quotient coordinate $(s-1/2)^2$ in Chapters~26 onward.\\
$F(w)$ & Unique entire quotient satisfying $\xi(1/2+z)=F(z^2)$.\\
$R_\star$ & Unique positive coefficient-majorant threshold satisfying $F(R_\star)=2F(0)$.\\
$\pi_n$ & Canonical half-mass coefficient law $a_nR_\star^n/(2F(0))$.\\
$R(s)$ & Candidate state readout $\braket{+}{\Psi(s)}$ in the information-spinor ansatz.\\
$g(s),h(s)$ & Non-vanishing scalar factors in local candidate factorizations; chapter-local declarations take precedence.\\
$\Phi(t)$ & Classical positive half-line Riemann Xi kernel in the normalization used by Chapters~27--55.\\
$K(t)$ & Even full-line Riemann kernel $K(t)=\Phi(|t|)/2$.\\
$L(t)$ & G024 log-potential $L(t)=-\log K(t)$.\\
$C(u)$ & G024 centred order-two correlation $\int_{\mathbb R}r^2K(u+r)K(u-r)\,dr$.\\
$D_y(u)$ & External Dimitrov--Xu tilt $\cosh(2yu)C(u)$.\\
$J_y(u)$ & Distinct internal Jensen tilt $\int r^2\cosh(2yr)K(u+r)K(u-r)\,dr$.\\
$H_y(q)$ & Radial-square external profile $D_y(\sqrt q)$ used for the complete-monotonicity route.\\
$\mu_u$ & Normalized G024 bridge probability measure with density proportional to $r^2K(u+r)K(u-r)$.\\
$A_u(r)$ & Bridge score observable $L'(u+r)+L'(u-r)$.\\
$B_u(r)$ & Bridge curvature observable $L''(u+r)+L''(u-r)$.\\
$C_u^{(3)}(r)$ & Third bridge observable $L'''(u+r)+L'''(u-r)$.\\
$R(u)$ & In G024, logarithmic correlation slope $-C'(u)/C(u)=\mathbb E_{\mu_u}[A_u]$.\\
$T_y(u)$ & External-tilt logarithmic slope $2|y|\tanh(2|y|u)$.\\
$N_y(u)$ & G024 surplus $R(u)-T_y(u)=-D_y'(u)/D_y(u)$.\\
$S_y(q)$ & Logarithmic complete-monotonicity slope $-\frac{d}{dq}\log H_y(q)=N_y(\sqrt q)/(2\sqrt q)$ for $q>0$.\\
$M_2(u)$ & Positive second-order radial margin $N_y+u(N_y^2-N_y')=4u^3H_y''(u^2)/H_y(u^2)$.\\
$F_m(q)$ & Normalized complete-monotonicity hierarchy $(-1)^mH_y^{(m)}(q)/H_y(q)$.\\
$\mu_3(A_u)$ & Third central moment of $A_u$ under the bridge measure.\\
$\widehat D_y$ & Fourier transform of the external correlation family; strict positivity for the full strip is the direct G024 open target.\\
\bottomrule
\end{longtable}

\section{Conventions}
All logarithms are natural. Entropies are measured in nats. Equality of state rays is denoted informally by $\sim$ or by square brackets $[\Psi]$. Numerical tolerances are not used in theorem statements. ``Zero'' without qualification means a non-trivial zeta zero when the context is the critical strip. Chapter-local symbol declarations override earlier local conventions. In particular, $u$, $q$, $R$, and $z$ are intentionally reused in distinct reductions; each chapter states its local definition before use.

\section{Proof-status vocabulary}
\begin{description}
\item[Canonical.] Promoted into the audited numbered theorem ledger through SOH-G023.
\item[Integrated research line.] Present on the repository mainline and included in the monograph, but not thereby promoted into the canonical numbered sequence; G024 currently has this status.
\item[Computer-assisted dependency.] A theorem step whose universal sign relies on a reproducible outward interval enclosure with an explicit analytic tail or equivalent trust boundary.
\item[Finite diagnostic.] Numerical evidence only; never a universal proof.
\item[OPEN.] A proof obligation not closed by the current repository state.
\end{description}
